\honest{Illusion of Transparency: Why No One Understands You}{Eliezer Yudkowsky}{2007}

``We always know what we mean by our words, and so we expect others to know it too.'' I call this the illusion of transparency, and I report several studies.

In the first, June recommends a restaurant to Mark. Mark leaves her a message saying that dinner there was ``marvelous, just marvelous.'' Of subjects told that the meal was bad, 59\% thought June would hear sarcasm. Of subjects told that it was good, 3\% did. Subjects told that the meal was bad but that Mark wanted to hide it thought June would not hear sarcasm in the same message. I call June ``Jane'' twice.

In another, two groups learned opposite meanings for the old idiom ``the goose hangs high.'' Each group thought an uninformed listener would pick the meaning it had been taught. In a footnote I list other idioms the study used and add, ``Ah, English, such a lovely language.''

In a third, speakers said ambiguous sentences such as ``The man is chasing a woman on a bicycle.'' They thought they were understood in 72\% of cases and were understood in 61\%. The same study had a second finding. People who overheard the speakers also knew the intended meaning, and they did not overestimate. So knowing what you mean is not enough to cause the illusion, which is the cause my first sentence gave. The authors put the bias in speakers judging ``their own utterances.'' I report the overhearers in one sentence and move on.

Then history. ``Two days before Germany's attack on Poland,'' Chamberlain sent Hitler a letter meant to say that Britain would fight, and Hitler heard it as conciliatory. The letter was sent ten days before the attack. It says the British government is ``resolved, and prepared, to employ without delay all the forces at their command.'' Hitler's reply restates the commitment and says, ``I take note of this statement.'' On the documents, he understood the letter and attacked anyway. My one example from outside the laboratory is of a message that was understood.

I end: ``Be not too quick to blame those who misunderstand your perfectly clear sentences.'' The sentences in the bicycle study were chosen for being ambiguous, so that study says nothing about how often clear ones are misread.
